\section{Generalized segment models and integrated likelihoods}
\label{sec:model}

\begin{figure*}[t]
  \centering
  \resizebox{0.94\textwidth}{!}{\begin{tikzpicture}[node distance=8mm and 10mm]
\node[primarybox,text width=31mm] (data) {Ordered observations\\$u,y,d,a$};
\node[secondarybox,right=of data,text width=35mm] (segment) {Observation-model calculation\\$\log A^{(r)}_{ij}$};
\node[primarybox,above right=7mm and 13mm of segment,text width=32mm] (sum) {Sum-product recursion\\marginal likelihoods and posterior marginals};
\node[primarybox,below right=7mm and 13mm of segment,text width=32mm] (max) {Max-sum recursion\\joint MAP partition};
\node[successbox,right=13mm of sum,text width=31mm] (post) {$p(k\mid y)$, changepoint probabilities, Bayes curves};
\node[successbox,right=13mm of max,text width=31mm] (map) {Joint MAP changepoints and fitted segments};
\draw[arrPrimary] (data)--(segment);
\draw[arrPrimary] (segment)--(sum);
\draw[arrPrimary] (segment)--(max);
\draw[arrPrimary] (sum)--(post);
\draw[arrPrimary] (max)--(map);
\node[draw=Rule,dashed,rounded corners,fit=(sum)(max),inner sep=4mm,label=above:{\small dynamic programming over ordered partitions}] {};
\end{tikzpicture}
}
  \caption{Statistical decomposition used by BayesBreak. A family-specific segment calculation supplies the integrated likelihood and optional posterior moment numerators. Local prior factors multiply the same triangular array. Sum-product and max-sum recursions over ordered partitions then produce posterior summaries and a joint MAP partition, respectively.}
  \label{fig:block-to-global}
\end{figure*}

\subsection{Ordered partitions and support}
Let $y_{1:n}$ be observed on an ordered design. Candidate boundary coordinates are denoted
$u_0<u_1<\cdots<u_n$; on a regular grid, $u_i=i$. A $k$-segment partition is a vector
$t=(t_0,\ldots,t_k)$ with
\[
\begin{aligned}
0=t_0<t_1<\cdots<t_k=n,\\
I_q(t)&=\{t_{q-1}+1,\ldots,t_q\}.
\end{aligned}
\]
The half-open notation $(i,j]$ denotes indices $i+1,\ldots,j$. Physical block length is
$\Delta_x(i,j)=u_j-u_i$ and is never inferred from an undefined observation $x_0$.

\begin{assumption}[Offline structural support and block independence]
\label{ass:offline}
The candidate segment-count set $\{1,\ldots,k_{\max}\}$ is finite. Structural admissibility of a segment is fixed by the candidate grid, declared length restrictions, and prior support before the response values are observed. A structurally admissible segment has finite nonnegative local evidence; zero evidence removes its likelihood contribution but does not alter the partition-prior normalizer. Structurally excluded segments have local prior factor zero and log-score $-\infty$. Conditional on a partition, segment parameters are independent across segments. The conditional partition prior factorizes through local terms as in Assumption~\ref{ass:factorized-prior}. For approximate segment routines, the same statements define an exact posterior for the supplied approximate score array; comparison to the intended nonconjugate model additionally requires Assumption~\ref{ass:block-error}.
\end{assumption}

\subsection{Exponential-family segment likelihoods with exposure and power weights}
Let $d_r$ collect scientifically specified family information, such as Poisson exposure, Binomial trial count, or Gaussian precision. Let $a_r\ge0$ be an optional power or reliability weight that multiplies the complete log contribution. The canonical local likelihood is
\begin{equation}
\begin{aligned}
\log p(y_r\mid\theta,d_r)^{a_r}
&=a_r\log h(y_r,d_r)\\
&\quad+a_rS(y_r,d_r)^\top\eta(\theta)\\
&\quad-a_rE(d_r)A(\theta).
\end{aligned}
\label{eq:local-ef}
\end{equation}
The corresponding block summaries are
\begin{align}
S_{ij}&=\sum_{r=i+1}^{j}a_rS(y_r,d_r),\nonumber\\
W_{ij}&=\sum_{r=i+1}^{j}a_rE(d_r),\nonumber\\
H_{ij}&=\sum_{r=i+1}^{j}a_r\log h(y_r,d_r).
\label{eq:block-summaries}
\end{align}
This notation prevents three distinct quantities from being collapsed into a single scalar ``weight'': design spacing affects the partition prior, $d_r$ changes the observation law, and $a_r$ raises the entire observation contribution to a power.

For conjugate hyperparameters $(\alpha,\beta)$, use the Diaconis--Ylvisaker normalizer convention
\begin{align}
p(\theta\mid\alpha,\beta)
&=\exp\{\eta(\theta)^\top\alpha-\beta A(\theta)-\log Z(\alpha,\beta)\},\nonumber\\
Z(\alpha,\beta)
&=\int e^{\eta(\theta)^\top\alpha-\beta A(\theta)}\dd\theta.
\label{eq:conjugate-prior}
\end{align}
Fix a scientifically declared reference descriptor $d_\star$ when the conditional observation mean depends on exposure or trial structure, and write $m_\star(\theta)=m(\theta;d_\star)$ for the segment functional reported by the analysis (for example, a Poisson rate at unit exposure rather than an exposure-specific count mean). Define
\begin{align*}
A^{(0)}_{ij}&=p(y_{i+1:j}\mid d_{i+1:j},a_{i+1:j}),\\
A^{(r)}_{ij}&=A^{(0)}_{ij}\,\mathbb E[m_\star(\theta)^r\mid y_{i+1:j}].
\end{align*}

\begin{theorem}[Exposure-aware conjugate block integration]
\label{thm:block-integration}
If $Z(\alpha_0+S_{ij},\beta_0+W_{ij})<\infty$, then
\begin{equation}
 A^{(0)}_{ij}=e^{H_{ij}}
 \frac{Z(\alpha_0+S_{ij},\beta_0+W_{ij})}{Z(\alpha_0,\beta_0)}.
\label{eq:block-evidence}
\end{equation}
The block posterior is conjugate with hyperparameters
$(\alpha_0+S_{ij},\beta_0+W_{ij})$. Whenever its $r$th moment exists,
\begin{equation}
 A^{(r)}_{ij}=A^{(0)}_{ij}
 \mathbb E[m_\star(\theta)^r\mid\alpha_0+S_{ij},\beta_0+W_{ij}].
\label{eq:block-moment}
\end{equation}
\end{theorem}
\begin{proof}
Multiplying the powered likelihood terms in Equation~\eqref{eq:local-ef} gives
$e^{H_{ij}}\exp\{\eta(\theta)^\top S_{ij}-W_{ij}A(\theta)\}$. Multiplication by the prior in Equation~\eqref{eq:conjugate-prior} produces the same conjugate kernel with updated hyperparameters $(\alpha_0+S_{ij},\beta_0+W_{ij})$. Integrating the kernel yields the ratio of normalizers in Equation~\eqref{eq:block-evidence}; division by this evidence gives the posterior. Inserting $m_\star(\theta)^r$ into the integral then factors the result into the evidence times the posterior moment, proving Equation~\eqref{eq:block-moment}.
\end{proof}

\begin{remark}[Scope of exactness]
Theorem~\ref{thm:block-integration} is exact for the displayed powered, exposure-aware model. It does not permit a Binomial trial count to be treated as a likelihood power, an irregular coordinate gap to be inserted as exposure, or a numerical integral to be described as conjugate closed form.
\end{remark}

\subsection{Local prior factors}
\begin{assumption}[Factorized ordered-partition prior]
\label{ass:factorized-prior}
For fixed $k$ and boundary coordinates $u$, the prior on $t$ is
\begin{equation}
 p(t\mid k,u)=C_k(u)^{-1}
 \prod_{q=1}^{k}c_u(t_{q-1},t_q)
 \prod_{q=1}^{k-1}h_u(t_q),
\label{eq:partition-prior}
\end{equation}
where $c_u(i,j)\ge0$ is a complete-segment cohesion and $h_u(j)\ge0$ is an interior-boundary hazard. Define
$\gamma_u(i,j)=c_u(i,j)h_u(j)^{\1\{j<n\}}$ and
$\widetilde A^{(r)}_{ij}=A^{(r)}_{ij}\gamma_u(i,j)$.
\end{assumption}
The normalizer $C_k(u)$ is itself a partition sum and can be computed by the same dynamic program after replacing all likelihood evidences by one. This ensures that fixed-$k$ evidence remains properly normalized for nonuniform local priors.

\begin{table*}[t]
\centering
\small
\caption{Supported local block constructions. Closed-form rows instantiate Theorem~\ref{thm:block-integration}; the Beta-valued response uses deterministic one-dimensional quadrature and is therefore exact only up to its numerical tolerance.}
\label{tab:block-families}
\begin{tabularx}{0.97\textwidth}{@{}L{0.15\textwidth}L{0.19\textwidth}L{0.17\textwidth}Y L{0.16\textwidth}@{}}
\toprule
Family & Descriptor $d_r$ & Segment parameter & Integrated block construction & Status \\
\midrule
Gaussian, known variance & precision or inverse variance & mean $\mu$ & Gaussian--Gaussian normalizer and posterior moments & closed form \\
Poisson & exposure $e_r$ & rate $\lambda$ & Gamma--Poisson normalizer; $E(d_r)=e_r$ & closed form \\
Binomial & trial count $m_r$ & probability $p$ & Beta--Binomial normalizer; $E(d_r)=m_r$ & closed form \\
Negative Binomial & fixed dispersion & success parameter & Beta-conjugate distributional parameterization & closed form \\
Beta-valued response & known precision $\phi_r$ & mean $\mu\in(0,1)$ & one-dimensional integral over a transformed mean & quadrature \\
Nonconjugate GLM & link-specific inputs & regression parameter & Laplace, variational, EP, PG, or quadrature block score & approximate unless certified \\
\bottomrule
\end{tabularx}
\end{table*}

\subsection{Segment marginal-likelihood interface}
The partition recursions require only a triangular array of segment log marginal likelihoods and, when posterior signal moments are requested, corresponding segment moment summaries. A family-specific implementation returns the segment log marginal likelihood, posterior hyperparameters or moments, a validity flag, and any numerical diagnostic needed to interpret an approximate value. Prefix summaries make Gaussian, Poisson, and Binomial segment calculations constant time after $\mathcal O(n)$ preprocessing. The recursion over partitions is therefore unchanged when the observation family changes.

\begin{specificationbox}[Required segment quantities]
For every candidate segment $(i,j]$, the family-specific calculation returns
$\bigl(\log A^{(0)}_{ij},\,\mathcal M_{ij},\,\texttt{status}_{ij},\,\texttt{diagnostics}_{ij}\bigr)$,
where $\mathcal M_{ij}$ contains the moments or posterior state required downstream. A segment outside the declared support receives log marginal likelihood $-\infty$, not an arbitrary finite penalty. Numerical approximations additionally report convergence information and a reference-error estimate whenever one is available.
\end{specificationbox}
